\section{Rank channels for retained vertex sets}\label{sec:env}

The hierarchy will use random retained sets to detect columns with many
substantial entries.  The atom construction will use the same sets to
represent matrix entries.  We define the sampling channels and prove the
marginal, conditional-pair, and count estimates needed for these two
constructions.  Section~\ref{sec:averaging} later constructs averaging
walks that preserve specified retention conditions.

The environment consists of the independent rank channels below together
with a finite string of independent uniform auxiliary bits.  The length
and uses of that string will be specified when needed.

We use a fixed number of independent rank channels, including channels
denoted $X,Y,Z$. Extra channels, when needed, are also a fixed finite
number. A channel is a uniformly random matrix
\[
A=(a_0,F)\in\mathbb F_P^{H\times h_0}
\]
conditioned on the last $r=h_0-1$ columns $F$ being linearly
independent. We call such an ordered collection of columns a \emph{frame}.
The integer $P$ is prime. It is chosen larger than every
identifier and sufficiently large for all rate comparisons below.
For a numerical identifier $y$, set
\[
v(y)=(1,y,\ldots,y^{h_0-1})^{\mathsf T},
\qquad a(y)=Av(y)\in\mathbb F_P^H.
\]
Field elements in a retention test use their representatives in
$\{0,\ldots,P-1\}$. At nominal rate $b_k$, let
\[
L_k=\lceil P2^{-k}\rceil,\qquad
\pi_{b_k}=(L_k/P)^H.
\]
The identifier is retained if all coordinates of $a(y)$ lie in
$\{0,\ldots,L_k-1\}$. Write $X_a(y)$ for the indicator on channel
$X$ at rate $a$, and similarly for other channels.
The same notation permits a specified fixed suffix in the channel
argument. Retention boxes are nested.
All rates in this section belong to the grid $b_k=2^{-Hk}$, and the
admissible range satisfies $P2^{-k}\ge1$. The choices of $H,h_0$ and
the identifier range are those of Section~\ref{sec:model}.

\begin{lemma}[Rank-channel estimates]\label{lem:env-ranks}
There are positive constants $c_\pi,C_\pi,C_{\rm pair},C_{\rm var}$,
depending only on $H,h_0$, with the following properties in every
admissible finite range.
\begin{enumerate}
\item A specified identifier is retained at rate $a$ with probability
$\pi_a$, where $c_\pi a\le\pi_a\le C_\pi a$.
\item For distinct identifiers $x,y$ and any admissible rates $u,a$,
\[
\Pr\{X_a(y)=1\mid X_u(x)=1\}\le C_{\rm pair}a.
\]
\item If $A_0$ is a deterministic set of $q$ distinct identifiers and
$T_a=\sum_{y\in A_0}X_a(y)$, then
\[
\Pr\{T_a<q\pi_a/2\}\le \frac{C_{\rm var}}{q\pi_a}
\quad(q>0).
\]
\item Replacing every channel argument $y$ by the same injective affine
map $\alpha y+\beta$ leaves the joint law of all evaluations unchanged,
provided $\alpha\ne0$ in $\mathbb F_P$. In particular, a common fixed
padding suffix leaves their law unchanged.
\end{enumerate}
\end{lemma}

\begin{proof}
The probability that an unrestricted $H$-by-$r$ matrix is a frame is
\[
q_{\rm fr}=\prod_{j=0}^{r-1}(1-P^{j-H})\ge c_{\rm fr}>0.
\]
For fixed positive-degree coefficients, $a_0$ is uniform, so any one
evaluation is uniform in $\mathbb F_P^H$ and independent of frame
success. This proves the exact marginal retention probability.
In the ranges used, take $P2^{-k}\ge1$. Then
$P2^{-k}\le L_k\le2P2^{-k}$, proving comparability. The prime
exponent can in fact make all these sides much larger than one.

Without conditioning on a frame, two evaluations at distinct field
points are independent uniform vectors: the two evaluation functionals
have rank two since $h_0\ge2$. Consequently
\[
\Pr\{X_a(y)=1,X_u(x)=1\mid F\text{ is a frame}\}
\le \frac{\pi_a\pi_u}{q_{\rm fr}}.
\]
The first marginal after conditioning is still $\pi_u$, which proves
the second assertion.

For the third assertion, in the unrestricted distribution the hit count
has mean $q\pi_a$ and variance at most $q\pi_a$, by pairwise
independence. Chebyshev's inequality bounds the stated lower-tail event
by $4/(q\pi_a)$. Conditioning on frame success increases this probability
by at most $1/q_{\rm fr}$.

Finally, substituting $\alpha y+\beta$ in a polynomial acts on its
positive-degree coefficient columns by an invertible triangular
matrix, with diagonal $\alpha,\ldots,\alpha^{h_0-1}$.
It therefore preserves the frame condition. Its action on the full
coefficient array is an invertible linear change of variables, so it
preserves uniform measure conditional on that condition. Appending
digits in radix $D_0$ has $\alpha$ a power of $D_0$; taking
$P>D_0$ ensures the required invertibility.
\end{proof}
